% ===================================
% W.A.S.T.E. Journal - LaTeX 论文模板
% ===================================
% 版本：1.0 (2026-03-08)
% 使用说明：
%   1. 使用 Overleaf 或本地 LaTeX 编辑器
%   2. 替换示例内容为你的论文
%   3. 编译生成 PDF
% ===================================

\documentclass[a4paper,12pt]{article}

% ===================================
% 基础包
% ===================================
\usepackage[utf8]{inputenc}
\usepackage[T1]{fontenc}
\usepackage{geometry}
\usepackage{graphicx}
\usepackage{hyperref}
\usepackage{amsmath}
\usepackage{amsfonts}
\usepackage{amssymb}
\usepackage{booktabs}
\usepackage{enumitem}
\usepackage{fancyhdr}
\usepackage{xcolor}
\usepackage{titlesec}
\usepackage{etoolbox}

% ===================================
% 页面设置
% ===================================
\geometry{
  top=2.5cm,
  bottom=2.5cm,
  left=2.5cm,
  right=2.5cm,
  headheight=15pt,
  footskip=1cm
}

% ===================================
% 颜色定义 (W.A.S.T.E. 主题色)
% ===================================
\definecolor{waste-charcoal}{RGB}{51,51,51}
\definecolor{waste-gold}{RGB}{197,160,47}
\definecolor{waste-red}{RGB}{209,47,47}
\definecolor{waste-gray}{RGB}{128,128,128}

% ===================================
% 页头页脚设置
% ===================================
\pagestyle{fancy}
\fancyhf{} % 清空所有页头页脚
\renewcommand{\headrulewidth}{2pt}
\renewcommand{\footrulewidth}{0.5pt}

% 页头：左侧 Logo + 期刊名，右侧栏目名
\fancyhead[L]{%
  \textbf{\Large\color{waste-charcoal} W.A.S.T.E.}\quad%
  \small\color{waste-gray}Journal%
}
\fancyhead[R]{%
  \small\color{waste-gold}\leftmark%
}

% 页脚：左侧作者，中间页码，右侧 DOI
\fancyfoot[L]{\small\color{waste-gray}\thepage}
\fancyfoot[C]{\small\color{waste-gray}W.A.S.T.E. Journal \textbullet\ Vol. 1, No. 1}
\fancyfoot[R]{\small\color{waste-gray}\texttt{doi:10.xxxx/waste.yyyy.xxxx}}

% 首页不显示页眉页脚
\fancypagestyle{plain}{%
  \fancyhf{}%
  \renewcommand{\headrulewidth}{0pt}%
  \renewcommand{\footrulewidth}{0pt}%
}

% ===================================
% 标题格式
% ===================================
\titleformat{\section}{\large\bfseries\color{waste-charcoal}}{\thesection}{1em}{}
\titleformat{\subsection}{\normalsize\bfseries\color{waste-charcoal}}{\subsection}{1em}{}
\titleformat{\subsubsection}{\normalsize\itshape\color{waste-charcoal}}{\subsubsection}{1em}{}

% ===================================
% 超链接设置
% ===================================
\hypersetup{
  colorlinks=true,
  linkcolor=waste-red,
  filecolor=waste-gold,
  urlcolor=waste-gold,
  citecolor=waste-red,
}

% ===================================
% 自定义命令
% ===================================
% 栏目命令
\newcommand{\wastesection}[1]{%
  \textcolor{waste-gold}{\textbf{#1}}%
}

% DOI 命令
\newcommand{\doi}[1]{%
  \texttt{\href{https://doi.org/#1}{doi:#1}}%
}

% 投稿日期命令
\newcommand{\submitted}[1]{%
  \small\color{waste-gray}Submitted: #1%
}

% ===================================
% 论文元数据（请修改此处）
% ===================================
\title{%
  \vspace{-2cm}%
  \textbf{\Huge\color{waste-charcoal} 你的论文标题}\\[0.5cm]
  \large\color{waste-gray}Your Paper Title in English%
}

\author{%
  \textbf{作者姓名}\thanks{Corresponding author: \texttt{email@example.com}}\\%
  \small 机构名称，城市，国家\\[0.3cm]%
  \textbf{Co-Author Name}\\%
  \small Another Institution, City, Country%
}

\date{%
  \wastesection{栏目名：填埋场 / Landfill}\\[0.2cm]%
  \submitted{2026-03-08}%
}

% ===================================
% 正文开始
% ===================================
\begin{document}

% 首页
\maketitle
\thispagestyle{plain} % 首页不使用页眉页脚

% ===================================
% 摘要
% ===================================
\begin{abstract}
\noindent 这里是摘要内容。摘要应该简洁明了地描述你的研究背景、方法、结果和结论。
对于 W.A.S.T.E. Journal，我们特别欢迎那些"失败"但具有教育意义的研究。

\vspace{0.5cm}
\noindent\textbf{关键词：} 关键词 1 \textbullet\ 关键词 2 \textbullet\ 关键词 3 \textbullet\ 关键词 4

\vspace{0.5cm}
\noindent\textbf{中图分类号：} XXXX \quad \textbf{文献标识码：} A
\end{abstract}

% ===================================
% 正文
% ===================================
\section{引言}

这里是引言部分。介绍研究背景、问题陈述和研究动机。

\subsection{研究背景}

详细说明研究背景，包括相关工作和现有研究的不足之处。

\subsection{问题陈述}

清晰描述你要解决的问题，以及为什么这个问题重要（或者为什么它不重要但有趣）。

\section{方法}

详细描述你的研究方法、实验设计和实施过程。

\subsection{实验设计}

\begin{itemize}[leftmargin=*]
  \item 实验环境：描述硬件/软件配置
  \item 数据集：说明数据来源和规模
  \item 评估指标：列出使用的评估标准
\end{itemize}

\subsection{实施细节}

提供足够的细节以便其他人可以复现你的"失败"。

\section{结果与讨论}

\subsection{实验结果}

展示你的结果，即使它们是阴性的或"失败"的。

\begin{table}[h]
\centering
\caption{实验结果对比}
\label{tab:results}
\begin{tabular}{lcc}
\toprule
方法 & 准确率 (\%) & 备注 \\
\midrule
Baseline & 85.2 & 参考文献 [1] \\
Our Method & 42.1 & 完全失败 \\
\bottomrule
\end{tabular}
\end{table}

\subsection{失败分析}

\textbf{这是 W.A.S.T.E. Journal 最看重的部分！} 详细分析为什么失败了，从中学到了什么。

\section{结论}

总结你的研究发现，即使结论是"此路不通"。

\section*{致谢}

感谢所有帮助过你的人，包括那些试图帮你但也没成功的。

% ===================================
% 参考文献
% ===================================
\begin{thebibliography}{99}

\bibitem{ref1}
Author, A. (2020).
\textit{Title of the Paper}.
Journal Name, 1(1), 1-10.

\bibitem{ref2}
Author, B. and Author, C. (2019).
\textit{Another Important Work}.
Conference Proceedings, pp. 100-200.

\end{thebibliography}

% ===================================
% 附录（可选）
% ===================================
\appendix
\section{附录 A：补充材料}

这里可以放置额外的数据、代码片段或详细推导。

% ===================================
% 作者贡献声明
% ===================================
\section*{作者贡献}

\textbf{作者 A：} 概念设计、实验实施、论文撰写 \\
\textbf{作者 B：} 数据分析、图表制作、论文修订

% ===================================
% 利益冲突声明
% ===================================
\section*{利益冲突声明}

作者声明无利益冲突。

% ===================================
% 数据可用性声明
% ===================================
\section*{数据可用性声明}

本研究的数据和代码已在 GitHub 公开：\url{https://github.com/your-repo}

% ===================================
% 许可证声明
% ===================================
\vspace{2cm}
\begin{center}
\small\color{waste-gray}
\textcopyright\ 2026 作者。本文采用 CC BY-NC-SA 4.0 许可证发布。\\
W.A.S.T.E. Journal \textbullet\ Worldwide Archive of Substandard Tech \& Experiments
\end{center}

\end{document}
